% Clase 8 --- Potencial de una barra en un punto del plano bisector (§6).
% Igual geometría que el cálculo del CAMPO en Clase 2, pero acá no hay que
% proyectar: se suma lambda dz / sqrt(y^2+z^2) y listo. La simetría sólo se usa
% para integrar de 0 a L/2 y duplicar.
\begin{center}
\begin{tikzpicture}[scale=1]

  % barra sobre el eje z (vertical)
  \draw[figline, line width=1.7pt] (0,-2.1) -- (0,2.1);
  \foreach \y in {-1.85,-1.4,...,2.0} {\node[figlbl,figred,scale=0.65] at (-0.22,\y){$+$};}
  \node[figlbl, figblue, anchor=south] at (0,2.2) {$\lambda$};
  \node[figlblaux, anchor=east] at (-0.45,2.1) {$L/2$};
  \node[figlblaux, anchor=east] at (-0.45,-2.1) {$-L/2$};

  % elemento dz
  \draw[figred, line width=2.4pt] (0,1.05) -- (0,1.38);
  \node[figlbl, figred, anchor=west] at (0.12,1.72) {$dQ = \lambda\,dz$};
  \node[figlblaux, anchor=east] at (-0.45,1.2) {$z$};
  \draw[figaux, line width=0.5pt] (0,1.2) -- (-0.4,1.2);

  % punto P
  \node[figneutra, minimum size=5pt] (P) at (3.0,0) {};
  \node[figlbl, anchor=north] at (3.0,-0.2) {$P$};
  \draw[figvecaux] (0,0) -- (2.85,0);
  \node[figlblaux, anchor=south] at (1.5,0.05) {$y$};

  % distancia del elemento a P
  \draw[figgray, line width=0.7pt] (0,1.2) -- (P);
  \node[figlblaux, anchor=south west] at (1.5,0.7) {$\sqrt{y^{2}+z^{2}}$};

  % el elemento simétrico
  \draw[figred, line width=2.4pt] (0,-1.38) -- (0,-1.05);
  \draw[figgray, line width=0.7pt, dashed] (0,-1.2) -- (P);
  \node[figlblaux, anchor=north] at (1.9,-1.15) {mismo aporte};

  % lectura
  \node[figlblaux, align=left, anchor=west] at (3.6,0.0)
    {sin proyectar: los dos aportes\\ simétricos se \textbf{suman} igual};

\end{tikzpicture}\\[0.5em]
{\small\itshape Compárese con el mismo problema para el \textbf{campo} (Clase
2): allí había que proyectar cada $d\vec E$ con un $\cos\theta$ y descartar la
componente que se cancelaba. Acá los dos elementos simétricos aportan
\textbf{lo mismo} y simplemente se suman, así que la integral se hace de $0$ a
$L/2$ y se duplica. La sustitución $z = y\sinh w$ la resuelve, y el resultado
sale en términos de $\operatorname{arcsinh}$.}
\end{center}
